%% abstract.tex

\chapter*{Abstract}

This dissertation presents the design, implementation, and evaluation of a
routing mechanism for a Mobile Ad-hoc Network (MANET) built on top of the
RA-TDMAs+ Time Division Multiple Access protocol. The work targets a practical
scenario in which a mobile robot streams video to a remote base station via an
intermediate relay node, operating without fixed network infrastructure.

The proposed system introduces a link quality metric derived from TDMA beacon
reception statistics and uses it to compute a Minimum Spanning Tree over the
network topology, from which per-destination next-hop decisions are derived.
Routes are installed directly into the Linux kernel routing table via Netlink,
and relay forwarding is implemented through a dual routing table architecture
that leverages the kernel's built-in IP forwarding without requiring a
user-space forwarding plane. A video quality feedback loop allows the base
station to signal stream loss to the robot node, triggering immediate relay
activation independently of the TDMA topology convergence cycle.

Experimental evaluation in a three-node hardware setup demonstrates routing
convergence after link failure, measures end-to-end round-trip latency under
direct and relay conditions, and characterises video stream delivery regularity
and throughput. Results confirm that the relay mechanism is transparent to the
video stream in terms of delivery rhythm, with the dominant convergence
bottleneck being the video loss detection timeout at the base station.
